\documentclass[10pt,a4paper]{amsart}
\usepackage[margin=19mm]{geometry}
\usepackage{amsmath,amssymb,mathtools}
\usepackage[scr=rsfs,cal=euler]{mathalfa}
\usepackage{xfrac}
\usepackage[unicode,hidelinks,bookmarks=false]{hyperref}
\newcommand{\ie}{i.e.\ }
\newcommand{\cf}{cf.\ }
\newcommand{\resp}{respectively\ }
\newcommand{\Subsection}{\subsection}
\providecommand{\nobreakdash}{\nobreak\hbox{-}}
\setlength{\emergencystretch}{2em}
\multlinegap0pt
\usepackage{amssymb}
\usepackage{enumerate}
\usepackage{cancel}
\newtheorem{lemma}{Lemma}
\newtheorem{thm}{Theorem}
\newcommand{\Hom}{{\on{Hom}}}
\newcommand{\Irr}{\operatorname{Irr}}
\newcommand{\on}{\operatorname}
\title{Admissibility of Bernstein centers}
\author{Tsao-Hsien Chen, Cheng-Chiang Tsai}
\date{Source: arXiv:2608.21981v1 (CC BY 4.0)}
\begin{document}
\maketitle

\noindent\emph{Source and licence.} Admissibility of Bernstein centers. Version arXiv:2608.21981v1 (CC BY 4.0), distributed by arXiv under the Creative Commons Attribution 4.0 International licence, http://creativecommons.org/licenses/by/4.0/, as recorded in the arXiv licence metadata for this exact version and shown on the abstract page https://arxiv.org/abs/2608.21981v1. Full licence notice: \texttt{licenses/arxiv-2608.21981-CC-BY-4.0.txt}.\par\smallskip

\noindent\textit{Editorial scope.} Counted unit: the article's thm main (source0.tex lines 484--491). The article's complete original proof of it is delivered with it (source0.tex lines 492--521, one \texttt{proof} environment of 30 lines). No mathematics is written, repaired or re-derived: the statement and the proof are the article's own lines, in the article's own notation, and no retained line is substituted or re-flowed.  Retained uncounted as the source-local context the proof needs: lines 435--453.  Nothing was omitted from any retained span, so the package carries no omission record.  No retained line contains a citation, so the wrapper carries no bibliography and no bibliography entry.  No retained line contains a figure environment, an image-inclusion command or drawing code, so no image is delivered and the package declares no asset dependency.  Editorial formatting only, in addition: an AMS \texttt{amsart} preamble that restates the article's own declarations and macros used by the retained lines, namely the environments \texttt{lemma}, \texttt{thm} on separate counters, together with the standard packages those lines need.  The retained environments print in this document as Lemma 1, Theorem 1 in order of appearance, whereas the article prints the same items with the numbering of its own full document.  The complete native source of the article as distributed in the arXiv e-print for this exact version is bundled unchanged as \texttt{sources/208290/source0.tex}.
    \begin{lemma}\label{projector}
Let $z\in Z(G)$, and let $(\rho,V_\rho)\in\Irr(K)$ be an irreducible representation of a compact open pro-$p$ subgroup $K\subset G$. Then
\[
f_z*\Theta_\rho\neq 0
\]
implies that there exists $\pi\in R(G)$ such that
\[
z(\pi)\neq 0
\qquad\text{and}\qquad
\operatorname{Hom}_K(\check\rho,\pi)\neq 0.
\]
%Here $\check\rho$ is the dual of $\rho$.
Moreover, if
\[
R(G)=\mathcal R(G)^{\heartsuit}\oplus
 R(G)^{\spadesuit}
\]
is a decomposition such that $z\in Z(G)^{\heartsuit}$, then $\pi$ may be chosen in $ R(G)^{\heartsuit}$.
\end{lemma}

\begin{thm}\label{main}
Suppose that $U\subset G$ is a compact open pro-$p$ subgroup and $R(G)=R(G)^{\heartsuit}\oplus R(G)^{\spadesuit}$ is a
 decomposition
%with the property that there exists  a compact open pro-$p$ subgroup $U\subset G$ such that 
such that any $(\pi,V_\pi)\in R(G)^\heartsuit$ is generated by its $U$-fixed vectors $(V_\pi)^U$. Then for any $z\in Z(G)^\heartsuit$,  the corresponding 
invariant distribution $f_z$ is $U$-admissible.
 
\end{thm}

\begin{proof}
Let $\rho\in\Irr(U')$ be an irreducible representation of a compact open subgroup $U'\subset U$
such that  $f_z*\Theta_\rho\neq0$.
By Lemma \ref{projector}, there exists 
 $(\pi,V_\pi)\in R(G)^\heartsuit$ 
 satisfying $z(\pi)\neq0$ and 
 $\Hom_{U'}(\check\rho,\pi)\neq 0$.
Choose a non-zero $U'$-equivariant embedding 
$f:V_{\check\rho}\hookrightarrow V_\pi$ and let 
$0\neq w$ lie in its image. 
Since $V_\pi$ is generated by its $U$-fixed vectors, we can 
write
$w=\sum_{j=1}^n g_jv_j$ for some
$v_j\in (V_\pi)^{U}$ and $g_j\in G$. Note that we have 
$g_jv_j\in V_{\pi}^{U'\cap g_jUg_j^{-1}}$ for 
any compact open $U'\subset U$ and 
$j=1,...,n$. On the other hand, we have
\[w=\pi(\Theta_\rho)(w)=\sum_{j=1}^n\pi(\Theta_\rho)(g_jv_j)\neq 0\]
and it follows that 
\[\pi(\Theta_\rho)(g_iv_i)\neq 0\]
for some $i\in [1,n]$.
Since  $\pi(\Theta_\rho):V_\pi\to \Hom_{U'}(\check\rho,\pi)\otimes V_{\check\rho}\subset V_\pi$
is $U'$-equivariant and 
$g_iv_i$ is fixed by $U'\cap g_iUg_i^{-1}$ we conclude that 
\[0\neq\pi(\Theta_\rho)(g_iv_i)\in\Hom_{U'}(\check\rho,\pi)\otimes (V_{\check\rho})^{U'\cap g_i Ug^{-1}}.\]
Thus $(V_{\check\rho})^{U'\cap g_i Ug^{-1}}\neq0$
and, since $U'\cap g_i Ug^{-1}$ is compact
 pro-$p$ and $\ell\neq p$, it implies that 
$(V_\rho)^{U'\cap g_i Ug^{-1}}\neq0$.
\end{proof}

\end{document}
